%%%%%%%%%%%%%%%%%
% Master Global CV - Guillaume Boileau - All skills consolidated
%%%%%%%%%%%%%%%%%

\documentclass[8pt,a4paper,ragged2e,withhyper]{../_shared/altacv}

\geometry{left=1.25cm,right=1.25cm,top=.5cm,bottom=1.cm,columnsep=1.2cm}

\usepackage{paracol}
\usepackage{fontawesome5}
\usepackage{hyperref}

\iftutex 
  \setmainfont{Roboto Slab}
  \setsansfont{Lato}
  \renewcommand{\familydefault}{\sfdefault}
\else
  \usepackage[rm]{roboto}
  \usepackage[defaultsans]{lato}
  \renewcommand{\familydefault}{\sfdefault}
\fi

\definecolor{SlateGrey}{HTML}{2E2E2E}
\definecolor{LightGrey}{HTML}{666666}
\definecolor{DarkPastelRed}{HTML}{8AC0D9}
\definecolor{PastelRed}{HTML}{00B0DA}
\definecolor{GoldenEarth}{HTML}{B4AD0C}

\colorlet{name}{black}
\colorlet{tagline}{PastelRed}
\colorlet{heading}{DarkPastelRed}
\colorlet{headingrule}{GoldenEarth}
\colorlet{subheading}{PastelRed}
\colorlet{accent}{PastelRed}
\colorlet{emphasis}{SlateGrey}
\colorlet{body}{LightGrey}

\definecolor{violet}{HTML}{5A2582}
\definecolor{rouge}{HTML}{F53B3B}
\definecolor{noir}{HTML}{00232C}
\definecolor{bleu}{HTML}{00B0DA}
\definecolor{rose}{HTML}{F831A0}
\definecolor{jaune}{HTML}{FFEC00}
\definecolor{vert}{HTML}{34AD6C}

\renewcommand{\namefont}{\Huge\rmfamily\bfseries}
\renewcommand{\personalinfofont}{\footnotesize}
\renewcommand{\cvsectionfont}{\LARGE\rmfamily\bfseries}
\renewcommand{\cvsubsectionfont}{\large\bfseries}
\renewcommand{\cvItemMarker}{{\small\textbullet}}
\renewcommand{\cvRatingMarker}{\faCircle}

\input{../_shared/pubs-num.tex}
\addbibresource{../_shared/sample.bib}

\begin{document}

\name{BOILEAU Guillaume}
\tagline{Ingénieur docteur | Data, logiciel, systèmes complexes, IVVQ, validation \& coordination technique}

\photoL{2.8cm}{../_shared/moi}

\personalinfo{%
  \email{guillaume.boileaupro@gmail.com}
  \phone{+33 6 59 94 85 84}
  \location{Alpes-Maritimes, FRANCE}
  \linkedin{guillaume-boileau-phd-891b81265}
  \researchgate{Guillaume-Boileau-2}
  \orcid{0000-0002-3576-6968}
}

\makecvheader
\vspace{-0.3cm}

\AtBeginEnvironment{itemize}{\small}
\columnratio{0.64}

\begin{paracol}{2}

\cvsection{Experience}

\cvevent{\textbf{Ingénieur logiciel full stack -- CPMS, intégration \& flux opérationnels}}{VEV Platform Services France}{2025 -- 2026}{Valbonne, France}

\textbf{Contexte :} Développement logiciel et support opérationnel pour une plateforme CPMS liée aux infrastructures de recharge de véhicules électriques.

\textbf{Objectif :} Contribuer au développement, à l'intégration et au suivi technique de services logiciels connectés à des équipements terrain et à des flux de données opérationnels.

\begin{itemize}
  \item \textbf{Développement logiciel :} Contribution au développement d'applications web et de services backend en TypeScript, Angular et Node.js.
  \item \textbf{Intégration système :} Travail sur des interfaces entre services logiciels, données opérationnelles et infrastructures de recharge.
  \item \textbf{Flux de données :} Suivi et analyse de flux opérationnels, notamment via des échanges FTP utilisés pour le pilotage et le suivi technique.
  \item \textbf{Support production :} Analyse d'incidents, investigation de comportements applicatifs et contribution à la fiabilisation des services.
  \item \textbf{Qualité et cohérence :} Vérification de la cohérence entre données, services applicatifs, comportement terrain et attentes opérationnelles.
  \item \textbf{Organisation agile :} Utilisation de Zenhub pour le suivi des tâches, la priorisation et la coordination dans un environnement de développement itératif.
  \item \textbf{Communication technique :} Formalisation d'observations, d'analyses d'incidents et de points techniques pour faciliter la prise de décision.
\end{itemize}

\divider

\cvevent{\textbf{Ingénieur R\&D senior -- Data, logiciel, validation \& coordination}}{Laboratoire Lagrange, Observatoire de la Côte d'Azur, CNES}{2023 -- 2025}{Nice, France}

\textbf{Contexte :} Activités d'ingénierie et de R\&D pour la mission spatiale LISA ESA/NASA, avec des enjeux de simulation, de traitement de données, d'intégration de modèles et de validation technique.

\textbf{Objectif :} Développer des outils Python permettant de structurer, comparer, intégrer et valider des modèles et des résultats issus de contributions hétérogènes.

\begin{itemize}
  \item \textbf{Développement Python :} Conception et développement de workflows pour l'intégration, la comparaison et la validation de modèles.
  \item \textbf{Pipelines de données :} Mise en place de chaînes de traitement pour organiser les entrées, transformer les sorties et comparer les résultats.
  \item \textbf{Validation technique :} Définition de contrôles de cohérence, de règles de comparaison et de méthodes de vérification des résultats.
  \item \textbf{Intégration de modèles :} Consolidation de résultats produits par différents modèles dans un cadre commun d'analyse et de revue.
  \item \textbf{Traçabilité :} Structuration des configurations, sorties de simulation, résultats d'analyse et éléments de documentation.
  \item \textbf{Documentation :} Utilisation de Confluence pour organiser la documentation technique, le suivi de validation et le partage de connaissances.
  \item \textbf{Coordination :} Travail avec des contributeurs scientifiques, logiciels et institutionnels dans un environnement international.
  \item \textbf{Plateforme logicielle :} Développement de CATpyLISA, plateforme Python pour l'intégration, la comparaison, la validation et la revue technique de modèles. \textcolor{DarkPastelRed}{\href{https://gitlab.oca.eu/gboileau/lisa_l3_tools}{\bf GitLab:CATpyLISA}}
\end{itemize}

\divider

\cvevent{\textbf{Data Scientist -- Modélisation, bruit \& analyse de données}}{Universiteit Antwerpen}{2021 -- 2023}{Antwerp, Belgium}

\textbf{Contexte :} Analyse de données et études de performance pour des instruments de précision dans un environnement technique international.

\textbf{Objectif :} Développer des workflows robustes et reproductibles pour identifier, caractériser et réduire des sources affectant la qualité de mesure.

\begin{itemize}
  \item \textbf{Analyse de données :} Traitement, structuration et analyse de données complexes issues d'instruments de précision.
  \item \textbf{Modélisation statistique :} Développement de méthodes d'analyse et d'inférence bayésienne, notamment par chaînes MCMC.
  \item \textbf{Traitement du signal :} Identification et caractérisation de contributions de bruit dans des mesures complexes.
  \item \textbf{Pipelines Python :} Mise en place de workflows de traitement, filtrage, transformation et analyse statistique.
  \item \textbf{Validation de résultats :} Analyse de sensibilité, vérification de cohérence et évaluation des limites de performance.
  \item \textbf{Reporting technique :} Production de résultats traçables, de synthèses techniques et de supports pour des parties prenantes internationales.
\end{itemize}

\divider

\cvevent{\textbf{Ingénieur R\&D junior -- Signal, simulation \& analyse de données}}{ARTEMIS, Observatoire de la Côte d'Azur}{2018 -- 2021}{Nice, France}

\textbf{Contexte :} Activités de R\&D liées à l'analyse de signaux faibles, à la simulation numérique et à la préparation de missions spatiales et d'instruments de détection.

\textbf{Objectif :} Développer des méthodes d'analyse, automatiser des simulations et produire des résultats exploitables dans un environnement scientifique et technique exigeant.

\begin{itemize}
  \item \textbf{Simulation numérique :} Développement d'approches de simulation pour des environnements de signaux complexes et des jeux de données de grande taille.
  \item \textbf{Analyse de données :} Mise en place de méthodes pour analyser et séparer des signaux faibles et superposés.
  \item \textbf{Traitement du signal :} Études de caractérisation, de décomposition et de comparaison de signaux dans des contextes bruités.
  \item \textbf{Automatisation :} Développement de workflows reproductibles pour les analyses répétées, les comparaisons et les études de sensibilité.
  \item \textbf{Investigation technique :} Analyse de sources de bruit et réalisation de diagnostics par corrélations entre capteurs.
  \item \textbf{Documentation :} Formalisation des méthodes, hypothèses, résultats et conclusions pour revue collaborative.
\end{itemize}

\switchcolumn

\cvsection{Profile}

\textbf{Ingénieur docteur avec une expérience transversale en data, logiciel, systèmes complexes, IVVQ, validation de résultats, traitement du signal, automatisation et coordination technique.}

Mon parcours combine développement Python, analyse de données, simulation numérique, intégration de modèles, automatisation de workflows, investigation d'anomalies, documentation technique et coordination dans des environnements exigeants : spatial, instrumentation, plateformes industrielles, CPMS et systèmes logiciels connectés.

Je peux me positionner sur des rôles d'ingénieur systèmes, IVVQ, RAMS/Safety support, data/AI engineer, backend Python, chef de projet technique, PMO technique, data product analyst ou coordinateur documentation / R\&T.

\begin{itemize}
  \item Culture forte de la qualité : cohérence, traçabilité, reproductibilité, validation et documentation.
  \item Expérience en développement logiciel, data processing, analyse statistique, simulation et test/validation.
  \item Capacité à faire le lien entre besoins métier, contraintes techniques, données, logiciel et résultats exploitables.
  \item Habitude des environnements collaboratifs, internationaux, multidisciplinaires et orientés livraison.
\end{itemize}

\cvsection{Languages}

\cvskill{English}{4}
\divider
\cvskill{Spanish}{2}
\divider

\medskip

\cvsection{Education}

\cvevent{PhD in Physics}{Université Côte d'Azur}{2018 -- 2021}{Nice, France}

\divider

\cvevent{Selective Graduate Program in Fundamental Physics (Magistère)}{Université Paris-Saclay}{2016 -- 2018}{Orsay, France}

\divider

\cvevent{MSc in Astronomy and Astrophysics}{Observatoire de Paris / Université Paris-Saclay}{2017 -- 2018}{Paris, France}

\divider

\cvevent{BSc in Fundamental Physics}{Université Paris-Saclay}{2015 -- 2016}{Orsay, France}

\cvsection{Professional Training}

\cvevent{Supervised Machine Learning: Regression \& Classification}{DeepLearning.AI -- Andrew Ng}{2026}{Coursera}

\divider

\cvevent{Practical AI-assisted development}{Prompting, code review, debugging, incremental development}{2026}{Self-training}

\end{paracol}

\cvsection{Projects}

\cvevent{CATpyLISA -- Plateforme Python d'intégration, comparaison et validation}{Mission spatiale LISA ESA/NASA}{2023 -- 2025}{}

Développement d'une plateforme Python dédiée à l'intégration de modèles, à la structuration de données, à la comparaison de résultats, à la validation technique et à la revue de workflows liés à la mission LISA.

\textbf{Rôle :} développement Python, conception de pipelines, structuration de données, logique de validation, contrôles de cohérence, documentation technique et coordination avec plusieurs contributeurs.

\textbf{Compétences mobilisées :} Python, data processing, automatisation, modélisation, validation, documentation, GitLab, environnement collaboratif international.

\divider

\cvevent{Flux de données opérationnels pour infrastructures de recharge}{Plateforme CPMS -- VEV Platform Services France}{2025 -- 2026}{}

Contribution à des services logiciels liés aux échanges de données opérationnelles et au comportement d'infrastructures connectées de recharge de véhicules électriques.

\textbf{Rôle :} développement logiciel, suivi de flux de données, analyse d'incidents, diagnostic opérationnel, fiabilisation d'interfaces et coordination technique.

\textbf{Compétences mobilisées :} TypeScript, Angular, Node.js, FTP, API, intégration logiciel/système, support production, analyse d'anomalies, Zenhub.

\divider

\cvevent{Collaborations scientifiques et techniques internationales}{LISA, LIGO--Virgo--KAGRA, Einstein Telescope}{2018 -- now}{}

Contribution à des travaux collaboratifs en analyse de données, traitement du signal, simulation, caractérisation de bruit, validation de résultats et investigation technique dans des environnements internationaux.

\textbf{Rôle :} développement d'outils d'analyse, contribution à des workflows de traitement, études techniques, documentation et revue collaborative.

\textbf{Compétences mobilisées :} analyse de données, traitement du signal, modélisation, simulation, validation, coordination technique et communication scientifique.

\cvsection{Strengths}

\vspace{-.3cm}

\cvsubsection{Data, analyse, modélisation \& IA}

{\small
\cvtag{Data Analysis}
\cvtag{Data Processing}
\cvtag{Data Quality}
\cvtag{Data Cleaning}
\cvtag{Data Transformation}
\cvtag{Data Pipelines}
\cvtag{Data Modeling}
\cvtag{Data Structuring}
\cvtag{Data Integration}
\cvtag{ETL Logic}
\cvtag{Data Migration Logic}
\cvtag{Data Consolidation}
\cvtag{Data Consistency}
\cvtag{Feature Extraction}
\cvtag{Predictive Modeling}
\cvtag{Classification}
\cvtag{Model Evaluation}
\cvtag{Cross-validation}
\cvtag{Bias/Variance}
\cvtag{Hyperparameter Tuning}
\cvtag{Machine Learning}
\cvtag{Deep Learning Awareness}
\cvtag{PyTorch}
\cvtag{Scikit-Learn}
\cvtag{Statistical Modeling}
\cvtag{Bayesian Inference}
\cvtag{MCMC}
\cvtag{Uncertainty Analysis}
\cvtag{Simulation Workflows}
\cvtag{Synthetic Data}
\cvtag{Model Comparison}
\cvtag{Model Validation}
\cvtag{LLM Awareness}
\cvtag{Prompt Engineering}
\cvtag{RAG Concepts}
\cvtag{Vector Search Concepts}
\cvtag{Agentic AI Concepts}
\cvtag{AI-Assisted Coding}
\cvtag{Power BI}
}

\divider\smallskip
\vspace{-.3cm}

\cvsubsection{Traitement du signal, physique \& systèmes de mesure}

{\small
\cvtag{Signal Processing}
\cvtag{Noise Analysis}
\cvtag{Noise Modeling}
\cvtag{Low-SNR Signals}
\cvtag{Waveform Simulation}
\cvtag{Spectral Analysis}
\cvtag{Physical-System Modeling}
\cvtag{Precision Instrumentation}
\cvtag{Environmental Noise}
\cvtag{Performance Analysis}
\cvtag{Sensitivity Studies}
\cvtag{Sensor Correlation Diagnostics}
\cvtag{System-Level Analysis}
\cvtag{Scientific Computing}
\cvtag{Acoustics Logic}
\cvtag{Ultrasound-Adjacent Signal Processing}
}

\divider\smallskip
\vspace{-.3cm}

\cvsubsection{Logiciel, backend, web \& automatisation}

{\small
\cvtag{Python}
\cvtag{NumPy}
\cvtag{SciPy}
\cvtag{Pandas}
\cvtag{Matlab}
\cvtag{C}
\cvtag{TypeScript}
\cvtag{JavaScript}
\cvtag{Angular}
\cvtag{Node.js}
\cvtag{Backend Engineering}
\cvtag{API Design}
\cvtag{REST API}
\cvtag{Webhooks Logic}
\cvtag{Asynchronous Workflows}
\cvtag{SQL}
\cvtag{Relational Databases}
\cvtag{NoSQL Awareness}
\cvtag{MongoDB}
\cvtag{Graph-Oriented Reasoning}
\cvtag{Complex Data Structures}
\cvtag{Workflow Automation}
\cvtag{Bash}
\cvtag{Shell Scripting}
\cvtag{Linux}
\cvtag{Unix}
\cvtag{Git}
\cvtag{GitLab}
\cvtag{GitHub}
\cvtag{Docker}
\cvtag{CI/CD}
\cvtag{FTP}
\cvtag{Prod / Preprod}
\cvtag{Monitoring}
\cvtag{Debugging}
\cvtag{Code Quality}
\cvtag{Maintainability}
\cvtag{VS Code}
\cvtag{LaTeX}
}

\divider\smallskip
\vspace{-.3cm}

\cvsubsection{IVVQ, test, validation \& qualité}

{\small
\cvtag{Verification \& Validation}
\cvtag{System Verification}
\cvtag{System Validation}
\cvtag{Software IVVQ}
\cvtag{Integration Testing}
\cvtag{Test Execution}
\cvtag{Test Procedures}
\cvtag{Test Automation}
\cvtag{Test Automation Scripting}
\cvtag{Validation Strategy Support}
\cvtag{Validation Campaigns}
\cvtag{Acceptance Criteria}
\cvtag{Requirements Verification}
\cvtag{Requirements Traceability}
\cvtag{Non-Regression Mindset}
\cvtag{Non-Regression Testing}
\cvtag{V-Cycle Logic}
\cvtag{Quality Assurance}
\cvtag{Quality Follow-up}
\cvtag{Validation Evidence}
\cvtag{Configuration Control}
\cvtag{Version Control}
\cvtag{Traceability}
\cvtag{Reproducibility}
\cvtag{Documentation Quality}
\cvtag{ISO 9001 Logic}
\cvtag{Continuous Improvement}
\cvtag{Technical Review Preparation}
}

\divider\smallskip
\vspace{-.3cm}

\cvsubsection{Systèmes, intégration, RAMS \& safety}

{\small
\cvtag{Systems Engineering}
\cvtag{Complex Systems}
\cvtag{System Integration}
\cvtag{Software Integration}
\cvtag{Hardware / Software Interfaces}
\cvtag{Operational Data Flows}
\cvtag{Interface Monitoring}
\cvtag{Technical Diagnostics}
\cvtag{Anomaly Investigation}
\cvtag{Incident Analysis}
\cvtag{Root Cause Analysis}
\cvtag{Issue Resolution}
\cvtag{Operational Reliability}
\cvtag{RAMS-Oriented Analysis}
\cvtag{Reliability Analysis}
\cvtag{Availability Awareness}
\cvtag{Maintainability Awareness}
\cvtag{Safety-Oriented Engineering}
\cvtag{Safety-Critical Awareness}
\cvtag{Risk Analysis}
\cvtag{Residual Risk Assessment}
\cvtag{Criticality Analysis}
\cvtag{Failure Scenario Analysis}
\cvtag{Failure Mode Awareness}
\cvtag{Risk Mitigation Support}
\cvtag{System Robustness}
\cvtag{Compliance Mindset}
\cvtag{EGSE-Relevant Interfaces}
\cvtag{Embedded Systems Awareness}
}

\divider\smallskip
\vspace{-.3cm}

\cvsubsection{Projet, PMO, product \& delivery}

{\small
\cvtag{Technical Project Management}
\cvtag{PMO Support}
\cvtag{Project Coordination}
\cvtag{Project Delivery}
\cvtag{Application Delivery}
\cvtag{Digital Platforms}
\cvtag{Data Platforms}
\cvtag{Product Requirements}
\cvtag{User Requirements}
\cvtag{Business Needs Translation}
\cvtag{Requirements Analysis}
\cvtag{Requirements Synthesis}
\cvtag{Roadmap Planning}
\cvtag{Roadmap Logic}
\cvtag{Backlog Logic}
\cvtag{Release Planning Awareness}
\cvtag{Work Package Structuring}
\cvtag{Task Follow-Up}
\cvtag{Task Allocation}
\cvtag{Contributor Alignment}
\cvtag{Stakeholder Management}
\cvtag{Stakeholder Coordination}
\cvtag{Cross-Functional Coordination}
\cvtag{Team Coordination}
\cvtag{Decision Support}
\cvtag{Risk Management}
\cvtag{Priority Management}
\cvtag{Planning}
\cvtag{Reporting}
\cvtag{Process Improvement}
\cvtag{Change Management Support}
\cvtag{Delivery Support}
\cvtag{Client-Oriented Mindset}
\cvtag{Customer/User Awareness}
}

\divider\smallskip
\vspace{-.3cm}

\cvsubsection{Documentation, méthodes \& outils collaboratifs}

{\small
\cvtag{Technical Documentation}
\cvtag{Documentation Management}
\cvtag{Knowledge Capitalisation}
\cvtag{Technical Synthesis}
\cvtag{Synthesis Reports}
\cvtag{Technical Reporting}
\cvtag{Process Documentation}
\cvtag{Review Preparation}
\cvtag{Technical Watch}
\cvtag{Data Repositories}
\cvtag{Archiving Logic}
\cvtag{Technical Libraries}
\cvtag{MindMap Logic}
\cvtag{Confluence}
\cvtag{Jira}
\cvtag{Zenhub}
\cvtag{Agile Environment}
\cvtag{Agile Tooling}
\cvtag{Scrum/Kanban Awareness}
\cvtag{Cycle en V}
\cvtag{GitLab Issues}
\cvtag{GitHub Projects Awareness}
\cvtag{Artifactory Awareness}
\cvtag{PyPI Awareness}
\cvtag{AsciiDoc Awareness}
}

\divider\smallskip
\vspace{-.3cm}

\cvsubsection{Domaines applicatifs}

{\small
\cvtag{Space Systems}
\cvtag{Satellite Mission Context}
\cvtag{ESA/NASA Collaboration}
\cvtag{LISA Mission}
\cvtag{Large-Scale International Projects}
\cvtag{Mission-Level Performance}
\cvtag{Scientific Payload Analysis}
\cvtag{Space Mission Data Processing}
\cvtag{Industrial Software}
\cvtag{EV Charging Infrastructure}
\cvtag{CPMS Platform}
\cvtag{Electrical Equipment Interfaces}
\cvtag{Digital Transformation}
\cvtag{Data-Driven Workflows}
\cvtag{AI-Oriented Workflows}
\cvtag{CRM Integration Logic}
\cvtag{ERP Integration Logic}
\cvtag{HubSpot Awareness / Ramp-Up}
\cvtag{n8n Awareness / Ramp-Up}
\cvtag{Low-Code Ramp-Up}
}

\divider\smallskip
\vspace{-.3cm}

\cvsubsection{HPC, calcul \& environnements techniques}

{\small
\cvtag{Slurm}
\cvtag{HTCondor}
\cvtag{Condor}
\cvtag{Distributed Workflows}
\cvtag{Large Datasets}
\cvtag{Automation Pipelines}
\cvtag{Reproducible Analyses}
\cvtag{Technical Environments}
}

\divider\smallskip
\vspace{-.3cm}

\cvsubsection{Soft skills}

{\small
\cvtag{Autonomy}
\cvtag{Analytical Thinking}
\cvtag{Problem Solving}
\cvtag{Execution Rigor}
\cvtag{Structured Execution}
\cvtag{Attention to Detail}
\cvtag{Adaptability}
\cvtag{Responsibility}
\cvtag{Fast Learner}
\cvtag{Learning Agility}
\cvtag{Growth Mindset}
\cvtag{Communication}
\cvtag{Structured Communication}
\cvtag{Technical English}
\cvtag{International Teamwork}
\cvtag{Collaboration}
\cvtag{Leadership}
\cvtag{Proactiveness}
\cvtag{Curiosity}
\cvtag{Reliability}
\cvtag{Synthesis}
\cvtag{Methodical Approach}
}

\end{document}